% MCQ Questions for Chapter 02
\section{বহুনির্বাচনি প্রশ্ন ও উত্তরমালা}
এগুলো বিষয়ভিত্তিক অনুশীলন। প্রচলিত সংখ্যাগত শ্রেণি ও তার প্রাসঙ্গিক ব্যাখ্যা মূল আলোচনায় রয়েছে।


\subsection*{\textbf{\underline{যোগাযোগব্যবস্থা ও ব্যান্ডউইথ}}}

\begin{multicols}{2}
\begin{enumerate}[series=chTwoRevision]

\item\chTwoMCQBlock{ডেটা কমিউনিকেশনের মূল উপাদান কয়টি?
\par ক. ৩টি \quad খ. ৪টি \par গ. ৫টি \quad ঘ. ৬টি }

\item\chTwoMCQBlock{প্রেরক থেকে প্রাপকের কাছে ডেটা যাওয়ার পথ কোন উপাদান?
\par ক. প্রোটোকল \quad খ. মাধ্যম \par গ. বার্তা \quad ঘ. এনকোডার }

\item\chTwoMCQBlock{ডিজিটাল নেটওয়ার্কে ব্যান্ডউইথ সাধারণত কী বোঝায়?
\par ক. সর্বোচ্চ ডেটা হার বা বহনক্ষমতা \quad খ. ডেটার আকার \par গ. নেটওয়ার্কের দূরত্ব \quad ঘ. সংকেতের শক্তি }

\item\chTwoMCQBlock{bps-এর পূর্ণরূপ কী?
\par ক. Bits per second \quad খ. Bytes per second \par গ. Binary per second \quad ঘ. Bits per system }

\item\chTwoMCQBlock{প্রচলিত পাঠ্যশ্রেণিতে ন্যারো ব্যান্ডের গতিসীমা কত?
\par ক. ৪৫-৩০০ bps \quad খ. ৯৬০০ bps \par গ. ১ Mbps \quad ঘ. ২ Mbps }

\item\chTwoMCQBlock{প্রচলিত HSC পাঠ্যশ্রেণিতে ব্রডব্যান্ডের প্রারম্ভিক হার কত?
\par ক. ৫১২ Kbps \quad খ. ১ Mbps \par গ. ২ Mbps \quad ঘ. ১০ Mbps }

\end{enumerate}
\end{multicols}

\subsection*{\textbf{\underline{ডেটা ট্রান্সমিশন মেথড ও মোড}}}

\begin{multicols}{2}
\begin{enumerate}[resume=chTwoRevision]

\item\chTwoMCQBlock{প্রচলিত ক্যারেক্টারভিত্তিক অ্যাসিনক্রোনাস ট্রান্সমিশনে ডেটা কীভাবে পাঠানো হয়?
\par ক. পৃথক ক্যারেক্টার ফ্রেমে \quad খ. একটি ব্লকের সমন্বিত ফ্রেমে \par গ. শুধু অবিচ্ছিন্ন অডিওপ্রবাহে \quad ঘ. সীমানাহীন এলোমেলো বিটে }

\item\chTwoMCQBlock{প্রচলিত ব্লকভিত্তিক সিনক্রোনাস ট্রান্সমিশনে ডেটা কীভাবে সংগঠিত হয়?
\par ক. প্রতিটি ক্যারেক্টারে পৃথক স্টার্ট ও স্টপ বিট দিয়ে \quad খ. ব্লক আকারে \par গ. কোনো ফ্রেমসীমা ছাড়াই \quad ঘ. শুধু একটি ক্যারেক্টারেই }

\item\chTwoMCQBlock{সিমপ্লেক্স মোডে ডেটা কোন দিকে প্রবাহিত হয়?
\par ক. একমুখী \quad খ. দ্বিমুখী, পর্যায়ক্রমে \par গ. দ্বিমুখী, একই সময়ে \quad ঘ. দিক নির্দিষ্ট নয় }

\item\chTwoMCQBlock{টেলিভিশন সম্প্রচারের একমুখী সংকেত পাঠানো কোন মোডের উদাহরণ?
\par ক. Simplex \quad খ. Half Duplex \par গ. Full Duplex \quad ঘ. Multiplex }

\item\chTwoMCQBlock{হাফ-ডুপ্লেক্স মোডের উদাহরণ কোনটি?
\par ক. ওয়াকি-টকি \quad খ. টেলিফোন \par গ. টেলিভিশন সম্প্রচার \quad ঘ. রেডিও সম্প্রচার }

\item\chTwoMCQBlock{ফুল-ডুপ্লেক্স মোডে ডেটা কীভাবে প্রবাহিত হয়?
\par ক. একমুখী \quad খ. দ্বিমুখী, পর্যায়ক্রমে \par গ. দ্বিমুখী, একই সময়ে \quad ঘ. দিক নির্দিষ্ট নয় }

\end{enumerate}
\end{multicols}

\subsection*{\textbf{\underline{ট্রান্সমিশন মিডিয়া}}}

\begin{multicols}{2}
\begin{enumerate}[resume=chTwoRevision]

\item\chTwoMCQBlock{UTP কেবলের পূর্ণরূপ কী?
\par ক. Unshielded Twisted Pair \quad খ. Universal Twisted Pair \par গ. Uniform Twisted Pair \quad ঘ. United Twisted Pair }

\item\chTwoMCQBlock{ফাইবার অপটিক কেবলে ডেটা কোন ধরনের সংকেতে বহন করা হয়?
\par ক. আলো \quad খ. বিদ্যুৎ \par গ. চুম্বক \quad ঘ. শব্দ }

\item\chTwoMCQBlock{উপযুক্ত নকশা ও যন্ত্রে ফাইবার অপটিক সংযোগের সুবিধা কোনটি?
\par ক. উচ্চ ডেটা হার \quad খ. দীর্ঘ দূরত্বে সংযোগ \par গ. EMI-এর প্রভাবমুক্ত আলোকসংকেত \quad ঘ. সবগুলো }

\item\chTwoMCQBlock{Bluetooth দিয়ে ব্যক্তিগত যন্ত্র সংযুক্ত করলে সাধারণত কোন এলাকার নেটওয়ার্ক হয়?
\par ক. PAN \quad খ. LAN \par গ. MAN \quad ঘ. WAN }

\item\chTwoMCQBlock{Wi-Fi প্রযুক্তি কোন IEEE মানের পরিবারের সঙ্গে সম্পর্কিত?
\par ক. IEEE 802.11\quad খ. IEEE 802.16
\par গ. IEEE 802.3\quad ঘ. IEEE 802.15.1}

\item\chTwoMCQBlock{প্রচলিত ডিশ অ্যান্টেনাভিত্তিক স্যাটেলাইট সংযোগে উচ্চ ফ্রিকোয়েন্সির কোন রেডিও তরঙ্গ ব্যবহৃত হয়?
\par ক. কম ফ্রিকোয়েন্সির রেডিও তরঙ্গ \quad খ. মাইক্রোওয়েভ \par গ. অবলোহিত রশ্মি \quad ঘ. দৃশ্যমান আলো }

\end{enumerate}
\end{multicols}

\subsection*{\textbf{\underline{মোবাইল যোগাযোগ ও নেটওয়ার্ক}}}

\begin{multicols}{2}
\begin{enumerate}[resume=chTwoRevision]

\item\chTwoMCQBlock{ডিজিটাল সেলুলার ডেটা ও SMS সেবার সূচনা কোন প্রজন্মের পরিচিত বৈশিষ্ট্য?
\par ক. 1G \quad খ. 2G \par গ. 3G \quad ঘ. 4G }

\item\chTwoMCQBlock{4G-এর প্রধান বৈশিষ্ট্য কোনটি?
\par ক. IP-ভিত্তিক নেটওয়ার্ক \quad খ. অ্যানালগ সংকেত \par গ. শুধু ভয়েস কল \quad ঘ. ডেটা যোগাযোগবিহীন ব্যবস্থা }

\item\chTwoMCQBlock{LAN-এর পূর্ণরূপ কী?
\par ক. Local Area Network \quad খ. Large Area Network \par গ. Long Area Network \quad ঘ. Limited Area Network }

\item\chTwoMCQBlock{প্রচলিত Ethernet LAN-এর স্টার বিন্যাসে কেন্দ্রীয় ডিভাইস হিসেবে সাধারণত কোনটি ব্যবহৃত হয়?
\par ক. Hub/Switch \quad খ. Router \par গ. Bridge \quad ঘ. Gateway }

\item\chTwoMCQBlock{ক্লাউড কম্পিউটিংয়ের একটি প্রধান সুবিধা কোনটি?
\par ক. প্রাথমিক হার্ডওয়্যার ব্যয় কমতে পারে এবং প্রয়োজনমতো সম্পদ বদলানো যায় \par খ. সব কাজের জটিল স্থানীয় সেটআপ বাধ্যতামূলক \par গ. শুধু স্থানীয় যন্ত্র থেকে ব্যবহার করা যায় \quad ঘ. সম্পদের পরিমাণ বদলানো যায় না }

\end{enumerate}
\end{multicols}

\subsection*{\textbf{\underline{উদ্দীপকভিত্তিক বহুনির্বাচনি প্রশ্ন}}}

\textbf{নিচের উদ্দীপকটি পড়ো এবং প্রশ্নের উত্তর দাও:}

\textit{একটি অফিসে ইন্টারনেট সংযোগের জন্য এমন একটি কেবল ব্যবহার করা হয় যা আলোর পূর্ণ অভ্যন্তরীণ প্রতিফলন পদ্ধতিতে ডেটা পাঠায়।}

\begin{multicols}{2}\begin{enumerate}[resume=chTwoRevision]\item\chTwoMCQBlock{উদ্দীপকে কোন কেবল ব্যবহৃত হয়েছে?
\par ক. টুইস্টেড পেয়ার \quad খ. কোঅ্যাক্সিয়াল \par গ. ফাইবার অপটিক \quad ঘ. তারবিহীন মাধ্যম }

\item\chTwoMCQBlock{উপযুক্ত সংযোগে এই কেবলের বৈশিষ্ট্য হলো—\par i. উচ্চ ডেটা হার সমর্থন করে; ii. আলোকসংকেত EMI-এর প্রভাবমুক্ত; iii. সংকেতের ক্ষয় তুলনামূলক কম।\par নিচের কোনটি সঠিক?
\par ক. i ও ii \quad খ. i ও iii \par গ. ii ও iii \quad ঘ. i, ii ও iii }

\end{enumerate}
\end{multicols}

\textbf{উত্তরমালা:}\par
1. গ, 2. খ, 3. ক, 4. ক, 5. ক, 6. খ, 7. ক, 8. খ, 9. ক।\par
10. ক, 11. ক, 12. গ, 13. ক, 14. ক, 15. ঘ, 16. ক, 17. ক, 18. খ।\par
19. খ, 20. ক, 21. ক, 22. ক, 23. ক, 24. গ, 25. ঘ।\par
